\documentclass[../BTL.tex]{subfiles}
\begin{document}
\section{Tập dữ liệu}
Tập dữ liệu là tập gồm các văn bản luật Việt Nam, các câu hỏi và danh sách các tài liệu liên quan cho câu hỏi đó. Dữ liệu này được đăng tải trên trang Huggingface \cite{dataset}, là một nền tảng chia sẻ tập dữ liệu machine learning phổ biến. Tập dữ liệu này được xây dựng nhằm mục đích phục vụ cho bài toán truy xuất thông tin pháp luật Việt Nam, một ứng dụng quan trọng trong việc tìm kiếm và tra cứu văn bản pháp luật một cách nhanh chóng và chính xác.

Các văn bản được đánh chỉ mục và xếp hạng dựa trên TF-IDF và BM25. Trong đánh giá, nhóm sử dụng kết hợp cả tập huấn luyện và tập kiểm thử, lấy 2000 mẫu đầu tiên (do giới hạn về phần cứng thiết bị) để chạy đánh giá và lấy kết quả. Việc giới hạn này là cần thiết do tập dữ liệu gốc có kích thước rất lớn, việc xử lý toàn bộ sẽ tốn nhiều thời gian và tài nguyên tính toán.

\subsection{Tập văn bản}
Tập văn bản (corpus) là tập hợp các văn bản pháp luật Việt Nam, được lưu trong file \texttt{corpus.csv} với 2 cột chính:
\begin{itemize}
    \item \textbf{cid}: ID duy nhất của mỗi văn bản, dùng để định danh và truy xuất.
    \item \textbf{text}: Nội dung đầy đủ của đoạn hoặc điều luật, là văn bản thô cần được xử lý.
\end{itemize}

Theo thống kê, có tổng cộng 222.981 tài liệu trong corpus, mỗi tài liệu là một đoạn văn bản pháp luật (có thể là một điều, khoản, điểm hoặc một đoạn văn bản trong các văn bản luật). Đây là một tập dữ liệu tương đối lớn, phản ánh khối lượng văn bản pháp luật khổng lồ của hệ thống pháp luật Việt Nam. Mỗi tài liệu thường có độ dài ngắn, tập trung vào một điều khoản cụ thể, điều này phù hợp cho việc truy xuất thông tin khi người dùng tìm kiếm một quy định cụ thể.

Đặc điểm của tập dữ liệu:
\begin{itemize}
    \item \textbf{Tính chuyên ngành cao:} Các văn bản sử dụng ngữ pháp và thuật ngữ pháp luật chuyên ngành.
    \item \textbf{Cấu trúc rõ ràng:} Các văn bản được phân chia theo điều, khoản, điểm.
    \item \textbf{Ngôn ngữ tiếng Việt:} Toàn bộ văn bản bằng tiếng Việt với đầy đủ dấu.
    \item \textbf{Độ dài ngắn:} Phần lớn các tài liệu có độ dài ngắn (trung bình 20-50 từ).
\end{itemize}

\subsection{Tập dữ liệu huấn luyện và tập dữ liệu kiểm thử}
Do đề tài sử dụng các phương pháp đánh giá dựa trên điểm (score-based) không phải huấn luyện mô hình máy học, do đó để đánh giá mô hình nhóm sử dụng kết hợp cả hai tập dữ liệu. Tuy nhiên do giới hạn phần cứng thiết bị, nhóm chỉ sử dụng 2000 mẫu đầu tiên để tiến hành đánh giá.

Tập dữ liệu có cấu trúc các cột như sau:

\begin{table}[H]
\centering
\begin{tabular}{|l|l|l|}
\hline
Cột & Kiểu & Mô tả \\ \hline
\texttt{qid} & int & ID câu hỏi duy nhất \\ \hline
\texttt{question} & string & Câu hỏi tìm kiếm \\ \hline
\texttt{cid} & string (list) & Danh sách ID văn bản liên quan \\ \hline
\texttt{context\_list} & string (list) & Nội dung đoạn văn liên quan \\ \hline
\end{tabular}
\end{table}

Tổng số câu hỏi của tập dữ liệu huấn luyện là 9.541 và tập dữ liệu kiểm thử là 3.184. Nhóm lấy 2.000 mẫu đầu tiên để làm nhiệm vụ đánh giá hệ thống truy xuất thông tin. Trong toàn bộ các tài liệu, mỗi truy vấn chỉ có 1 hoặc một vài tài liệu liên quan. Điều này cho thấy tỉ lệ mẫu dương tính (positive) rất thấp - đây là thách thức lớn cho các hệ thống truy xuất thông tin vì:

\begin{itemize}
    \item Hệ thống phải tìm một "kim đốt giữa mo" trong hàng trăm nghìn tài liệu.
    \item Precision@K sẽ bị ảnh hưởng nghiêm trọng khi K nhỏ.
    \item Recall@10 đạt ~77\% là kết quả rất tốt trong điều kiện này.
\end{itemize}

Đặc điểm của tập dữ liệu đánh giá:
\begin{itemize}
    \item \textbf{Tài liệu liên quan thưa thớt:} Mỗi truy vấn chỉ có vài tài liệu liên quan.
    \item Các câu hỏi được viết bằng ngôn ngữ tự nhiên.
    \item Bao gồm nhiều lĩnh vực pháp luật khác nhau.
    \item Có sẵn danh sách tài liệu liên quan làm ground truth.
\end{itemize}

Để hiểu rõ hơn về độ khó của bài toán, ta có thể phân tích tỷ lệ tài liệu liên quan trên tổng số tài liệu. Với tổng cộng 222.981 tài liệu trong corpus và trung bình chỉ có khoảng 1-3 tài liệu liên quan cho mỗi truy vấn, tỷ lệ này dao động từ 0.0004\% đến 0.0013\%. Đây là một tỷ lệ cực kỳ thấp, đặt ra thách thức lớn cho các thuật toán truy xuất thông tin. Hệ thống cần phải tìm chính xác những tài liệu thuộc phần nhỏ bé này trong hàng trăm nghìn tài liệu khác.

Một đặc điểm quan trọng khác của tập dữ liệu là sự đa dạng về chủ đề pháp lý. Các câu hỏi trong tập dữ liệu bao phủ nhiều lĩnh vực như luật hình sự, luật dân sự, luật hành chính, luật thương mại, luật đất đai, luật lao động, v.v. Điều này đảm bảo rằng kết quả đánh giá không bị thiên lệch về một lĩnh vực cụ thể mà phản ánh hiệu suất tổng thể của hệ thống trên nhiều loại truy vấn khác nhau.


\section{Xây dựng chỉ mục Elasticsearch}
Elasticsearch là một công cụ tìm kiếm phân tán mạnh mẽ, được sử dụng rộng rãi trong các ứng dụng thực tế. Trong đề tài này, Elasticsearch đóng vai trò là baseline để so sánh với các mô hình cài đặt thủ công. Việc sử dụng Elasticsearch làm baseline mang lại nhiều lợi ích: thứ nhất, đây là một hệ thống tìm kiếm đã được tối ưu hóa và chứng minh hiệu quả trong nhiều ứng dụng công nghiệp; thứ hai, Elasticsearch cung cấp nhiều tính năng mạnh mẽ như tìm kiếm full-text, phân tích ngôn ngữ tự nhiên, và khả năng mở rộng quy mô; thứ ba, việc so sánh với Elasticsearch giúp đánh giá xem các cài đặt thủ công có thể cạnh tranh được với các giải pháp sẵn có hay không.

Trong đề tài này, nhóm sử dụng phiên bản Elasticsearch 8.x với Python client tương ứng. Cấu hình cluster được thiết kế đơn giản phù hợp với môi trường thử nghiệm: một node master-dual được cấu hình với 2GB heap memory. Mặc dù đây là cấu hình nhỏ, nó đủ để xử lý tập dữ liệu 222.981 văn bản pháp lý một cách hiệu quả.

\subsection*{Hai phương pháp đánh chỉ mục}
Với bước xây dựng chỉ mục, đề tài sử dụng 2 cách để đánh giá và so sánh:

\begin{enumerate}
    \item \textbf{ES với chỉ mục mặc định (Normal index):} Sử dụng bộ phân tích (analyzer) mặc định của Elasticsearch, được thiết kế cho tiếng Anh. Bộ phân tích này sẽ tách từ theo dấu cách và dấu câu.
    \item \textbf{ES với chỉ mục đã xử lý (Processed index):} Tiến hành tách từ tiếng Việt bằng thư viện underthesea trước, sau đó nối lại theo dấu cách và đưa vào chỉ mục của Elasticsearch.
\end{enumerate}

Sự khác biệt giữa hai phương pháp:
\begin{itemize}
    \item \textbf{Chỉ mục bình thường:} Elasticsearch tự động tokenize và phân tích văn bản.
    \item \textbf{Chỉ mục đã xử lý:} Văn bản đã được underthesea xử lý trước khi đưa vào Elasticsearch.
\end{itemize}

Mục tiêu của việc so sánh này là chứng minh tầm quan trọng của tiền xử lý văn bản tiếng Việt trong hệ thống tìm kiếm. Khi sử dụng bộ phân tích mặc định của Elasticsearch, văn bản tiếng Việt sẽ bị tách từ một cách không chính xác do đặc điểm của tiếng Việt là không có dấu cách giữa các từ trong từ ghép. Ví dụ, câu "Người phạm tội sẽ bị truy cứu trách nhiệm hình sự" sẽ bị tách thành các token riêng lẻ như "Người", "phạm", "tội", "sẽ", "bị", "truy", "cứu", "trách", "nhiệm", "hình", "sự" thay vì các từ ghép có nghĩa như "Người", "phạm", "tội", "truy", "cứu", "trách", "nhiệm", "hình", "sự".

Điều này dẫn đến việc tìm kiếm với từ khóa "truy cứu" sẽ không khớp với văn bản vì Elasticsearch đã tách "truy" và "cứu" thành hai token riêng biệt. Ngược lại, khi văn bản đã được xử lý trước bằng underthesea, từ ghép "truy\_cứu" được giữ nguyên và việc tìm kiếm sẽ chính xác hơn.

Kết quả thực nghiệm cho thấy chỉ mục đã xử lý cho kết quả tốt hơn, chứng tỏ tiền xử lý văn bản tiếng Việt là cần thiết.

Quy trình đánh chỉ mục được thực hiện như sau:

\begin{enumerate}
    \item \textbf{Đọc dữ liệu:} Đọc file corpus.csv chứa 222.981 văn bản pháp lý.
    \item \textbf{Xử lý trước (cho Processed index):} Với mỗi văn bản, sử dụng underthesea để tách từ và nối từ ghép bằng dấu "\_".
    \item \textbf{Tạo index:} Tạo index trong Elasticsearch với mapping định nghĩa trước.
    \item \textbf{Index documents:} Đưa từng văn bản vào index với các trường cid và content.
    \item \textbf{Refresh:} Refresh index để các document có thể được tìm kiếm ngay lập tức.
\end{enumerate}

Thời gian đánh chỉ mục cho toàn bộ 222.981 văn bản là khoảng 5-10 phút tùy thuộc vào cấu hình phần cứng. Kích thước index sau khi đánh chỉ mục là khoảng 150-200 MB, bao gồm dữ liệu văn bản và các cấu trúc phục vụ tìm kiếm.

\subsection*{Cấu hình Mapping}
Mapping là cấu hình cấu trúc dữ liệu cho index trong Elasticsearch. Mapping định nghĩa cách dữ liệu được lưu trữ và index, bao gồm kiểu dữ liệu, cách phân tích, và các thuộc tính khác. Với trường content trong dữ liệu, kiểu dữ liệu khi ánh xạ vào chỉ mục của Elasticsearch là \texttt{text}.

\begin{lstlisting}[language=Python]
mapping = {
    "mappings": {
        "properties": {
            "cid": {"type": "keyword"},
            "content": {
                "type": "text",
                "analyzer": "standard",
                "fields": {
                    "keyword": {"type": "keyword"}
                }
            }
        }
    }
}
\end{lstlisting}

Elasticsearch sẽ phân tích, phân tách nội dung thành các token riêng lẻ và phù hợp cho việc tìm kiếm tài liệu chứa các từ. Tuy nhiên, trường kiểu \texttt{text} không thể dùng cho sắp xếp, tổng hợp hoặc tìm kiếm chính xác tuyệt đối (exact match).

Phần tiếp theo của cấu hình mapping là phần analyzer, sử dụng bộ phân tích chuẩn (standard analyzer). Bộ phân tích này của Elasticsearch gồm 3 thành phần:

\begin{enumerate}
    \item \textbf{Tokenizer:} Tách văn bản thành từng token theo dấu cách, dấu câu.
    \item \textbf{Lowercase filter:} Chuyển tất cả các ký tự về chữ thường.
    \item \textbf{Stop Token Filter:} Loại bỏ từ dừng tiếng Anh mặc định.
\end{enumerate}

Bộ phân tích standard được thiết kế chủ yếu cho tiếng Anh, do đó khi áp dụng cho tiếng Việt, nó chỉ có thể tách từ dựa trên dấu cách. Điều này không phù hợp với đặc điểm của tiếng Việt - ngôn ngữ sử dụng dấu cách để phân tách từ trong câu, nhưng các từ ghép không có dấu cách giữa các tiếng. Vì vậy, để cải thiện hiệu quả tìm kiếm cho tiếng Việt, cần sử dụng các phương pháp tiền xử lý như tách từ bằng underthesea trước khi đưa vào Elasticsearch.

Phần tiếp theo là trường keywords - tạo thêm một trường con \texttt{content.keyword} bên cạnh \texttt{content}. Bảng \ref{tab:content-vs-content-keyword} so sánh giữa 2 trường content và content.keyword.

\begin{table}[H]
\centering
\begin{tabular}{|l|l|l|}
\hline
Thuộc tính & content (text) & content.keyword (keyword) \\ \hline
Loại & Full-text & Exact match \\ \hline
Phân tích & Có (tokenize) & Không (giữ nguyên) \\ \hline
Dùng cho & Match query, search & Sắp xếp, tổng hợp, lọc \\ \hline
Ví dụ & Tìm một từ trong văn bản & So sánh chính xác chuỗi \\ \hline
\end{tabular}
\caption{So sánh trường \texttt{content} và trường \texttt{content.keyword}}
\label{tab:content-vs-content-keyword}
\end{table}

Khi phân tích tác động đến kết quả đánh giá, dễ thấy rằng với ngôn ngữ tiếng Việt, chỉ mục đã được tách từ sử dụng thư viện underthesea cho kết quả tốt hơn khi sử dụng cách tách từ chuẩn của Elasticsearch. Điều này chứng tỏ ngôn ngữ mặc định áp dụng của Elasticsearch là tiếng Anh không phù hợp với tiếng Việt, và việc tiền xử lý là cần thiết.

Để tối ưu hóa hiệu suất tìm kiếm, Elasticsearch còn hỗ trợ nhiều loại query khác nhau. Trong đề tài này, nhóm sử dụng \texttt{MultiMatchQuery} với trường \texttt{content} để thực hiện tìm kiếm full-text. Query này cho phép tìm kiếm một từ hoặc cụm từ trên nhiều trường cùng lúc và tính điểm dựa trên mức độ khớp.


\section{Xây dựng chỉ mục ngược}
\subsection{Quy trình xây dựng chỉ mục ngược}

Chỉ mục ngược được xây dựng từ tập dữ liệu văn bản pháp luật đã qua tiền xử lý. Quy trình xây dựng chỉ mục ngược bao gồm các bước sau:

\textbf{Bước 1: Tạo cặp (term, docID)}

Với mỗi tài liệu trong corpus, sau khi tách từ và loại bỏ từ dừng, tạo ra danh sách các cặp (term, docID). Ví dụ, nếu tài liệu Doc1 chứa các từ "thuế", "luật", "phạt", ta tạo ra các cặp: (thuế, Doc1), (luật, Doc1), (phạt, Doc1).

\textbf{Bước 2: Sắp xếp}

Sắp xếp tất cả các cặp (term, docID) theo thứ tự từ điển của term, sau đó theo docID. Độ phức tạp của bước này là $O(N \times L \times \log(N \times L))$, với $N$ là số tài liệu và $L$ là độ dài trung bình của tài liệu.

\textbf{Bước 3: Gộp postings}

Duyệt qua danh sách đã sắp xếp và gộp các mục trùng lặp. Với mỗi term, tạo postings list chứa danh sách các docID và tần suất xuất hiện. Cấu trúc dữ liệu được tổ chức dưới dạng từ điển với key là term, value chứa doc\_freq (số tài liệu chứa term) và postings list (danh sách các cặp docID và term frequency).

\textbf{Bước 4: Lưu trữ}

Lưu chỉ mục ngược vào file để tái sử dụng mà không cần xây dựng lại. Chỉ mục được lưu dưới dạng JSON để dễ đọc và xử lý. Kích thước chỉ mục khoảng 150-200 MB cho 222,981 tài liệu.

\subsection{Tối ưu hóa chỉ mục}

Để tăng tốc độ truy vấn, một số tối ưu hóa được áp dụng:

\begin{itemize}
    \item \textbf{Chỉ mục trong bộ nhớ:} Toàn bộ chỉ mục được nạp vào RAM khi khởi động hệ thống, giúp truy vấn nhanh mà không cần đọc từ đĩa.
    
    \item \textbf{Cấu trúc dữ liệu tối ưu:} Sử dụng dictionary của Python với độ phức tạp truy cập $O(1)$ cho việc tìm kiếm term.
    
    \item \textbf{Danh sách postings được sắp xếp:} Mỗi postings list được sắp xếp theo docID để tối ưu các phép toán tập hợp (AND, OR).
\end{itemize}

\subsection{Các phép toán trên chỉ mục ngược}

Các phép toán tập hợp trên chỉ mục ngược được thực hiện bằng thuật toán merge với độ phức tạp tuyến tính.

\textbf{Phép AND:} Tìm các tài liệu chứa cả hai từ.

\begin{algorithm}[H]
\caption{Intersect (Phép AND)}
\SetAlgoLined
\KwIn{$P_1, P_2$ - hai danh sách postings đã sắp xếp}
\KwOut{$R$ - danh sách kết quả}
$i \leftarrow 0$, $j \leftarrow 0$, $R \leftarrow []$\;
\While{$i < |P_1|$ \textbf{and} $j < |P_2|$}{
    \eIf{$P_1[i] = P_2[j]$}{
        Thêm $P_1[i]$ vào $R$\;
        $i \leftarrow i + 1$, $j \leftarrow j + 1$\;
    }{
        \eIf{$P_1[i] < P_2[j]$}{
            $i \leftarrow i + 1$\;
        }{
            $j \leftarrow j + 1$\;
        }
    }
}
\Return{$R$}
\end{algorithm}

\textbf{Phép OR:} Tìm các tài liệu chứa ít nhất một trong hai từ.

\begin{algorithm}[H]
\caption{Union (Phép OR)}
\SetAlgoLined
\KwIn{$P_1, P_2$ - hai danh sách postings đã sắp xếp}
\KwOut{$R$ - danh sách kết quả}
$i \leftarrow 0$, $j \leftarrow 0$, $R \leftarrow []$\;
\While{$i < |P_1|$ \textbf{and} $j < |P_2|$}{
    \eIf{$P_1[i] = P_2[j]$}{
        Thêm $P_1[i]$ vào $R$\;
        $i \leftarrow i + 1$, $j \leftarrow j + 1$\;
    }{
        \eIf{$P_1[i] < P_2[j]$}{
            Thêm $P_1[i]$ vào $R$, $i \leftarrow i + 1$\;
        }{
            Thêm $P_2[j]$ vào $R$, $j \leftarrow j + 1$\;
        }
    }
}
Nối phần còn lại của $P_1$ hoặc $P_2$ vào $R$\;
\Return{$R$}
\end{algorithm}

\textbf{Phép NOT:} Tìm các tài liệu không chứa từ.

\begin{algorithm}[H]
\caption{Complement (Phép NOT)}
\SetAlgoLined
\KwIn{$P$ - danh sách postings của từ cần loại, $D$ - tập tất cả docID}
\KwOut{$R$ - danh sách kết quả}
$R \leftarrow D$\;
\ForEach{$doc \in P$}{
    Loại bỏ $doc$ khỏi $R$\;
}
\Return{$R$}
\end{algorithm}


\section{Các phương pháp truy xuất}
\subsection{Mô hình Boolean}

Mô hình Boolean được cài đặt dựa trên các phép toán tập hợp trên chỉ mục ngược. Mỗi truy vấn chỉ chứa tối đa một toán tử (AND, OR, hoặc NOT) để đơn giản hóa việc xử lý.

\textbf{Cấu trúc truy vấn:}

Truy vấn Boolean được biểu diễn dưới dạng danh sách các token, trong đó có thể chứa một toán tử (AND, OR, NOT) để phân tách hai nhóm term:
\begin{itemize}
    \item \texttt{first\_terms}: Các term trước toán tử
    \item \texttt{operator}: Toán tử AND, OR, hoặc NOT
    \item \texttt{second\_terms}: Các term sau toán tử
\end{itemize}

Ví dụ: truy vấn "thuế AND luật" sẽ được tách thành first\_terms = ["thuế"], operator = "AND", second\_terms = ["luật"].

\textbf{Thuật toán tìm kiếm Boolean:}

\begin{algorithm}[H]
\caption{BooleanSearch}
\SetAlgoLined
\KwIn{$query$ - danh sách các token truy vấn}
\KwOut{$result$ - danh sách các docID}
$first\_terms \leftarrow []$, $second\_terms \leftarrow []$, $operator \leftarrow \text{""}$
$is\_first \leftarrow \text{true}$

\ForEach{$term \in query$}{
    \If{$term \in \{AND, OR, NOT\}$}{
        $is\_first \leftarrow \text{false}$
        $operator \leftarrow term$
        \textbf{continue}
    }
    \If{$is\_first$}{
        Thêm $term$ vào $first\_terms$
    }{
        Thêm $term$ vào $second\_terms$
    }
}

\If{$operator = \text{""}$}{
    $operator \leftarrow \text{AND}$\tcc*{Mặc định AND nếu không có toán tử}
}

\Switch{$operator$}{
    \Case{AND}{
        $all\_terms \leftarrow first\_terms + second\_terms$
        $result \leftarrow \text{getDocsContaining}(all\_terms[0])$
        \ForEach{$term \in all\_terms[1:]$}{
            $result \leftarrow result \cap \text{getDocsContaining}(term)$
        }
    }
    \Case{OR}{
        $result \leftarrow \emptyset$
        \ForEach{$term \in first\_terms$}{
            $result \leftarrow result \cup \text{getDocsContaining}(term)$
        }
        \ForEach{$term \in second\_terms$}{
            $result \leftarrow result \cup \text{getDocsContaining}(term)$
        }
    }
    \Case{NOT}{
        $all\_docs \leftarrow \text{getAllDocIds}()$
        $excluded \leftarrow \emptyset$
        \ForEach{$term \in first\_terms + second\_terms$}{
            $excluded \leftarrow excluded \cup \text{getDocsContaining}(term)$
        }
        $result \leftarrow all\_docs - excluded$
    }
}
\Return{$result$}
\end{algorithm}

\textbf{Hạn chế:} Do không có cơ chế xếp hạng, tất cả tài liệu khớp được trả về với cùng mức độ ưu tiên.


\subsection{Mô hình không gian vector (VSM)}

VSM được cài đặt với các bước sau:

\textbf{Bước 1: Xây dựng không gian vector}

Xác định không gian vector với số chiều bằng số từ duy nhất trong toàn bộ corpus (khoảng 50,000 chiều sau khi loại bỏ từ dừng). Mỗi tài liệu và truy vấn được biểu diễn dưới dạng vector trong không gian này.

\textbf{Bước 2: Tính trọng số TF-IDF}

Với mỗi từ $t$ trong tài liệu $d$, trọng số TF-IDF được tính như sau:

Tần suất xuất hiện của từ $t$ trong tài liệu $d$ được tính bằng số lần xuất hiện của $t$ chia cho tổng số từ trong $d$. Để giảm ảnh hưởng của các tài liệu dài, có thể sử dụng công thức logarit: $TF(t, d) = 1 + \log(\text{số lần xuất hiện của } t)$.

IDF của từ $t$ được tính bằng logarit của tổng số tài liệu chia cho số tài liệu chứa $t$, với công thức smoothing để tránh chia cho 0: $IDF(t) = \log\left(\frac{N + 1}{n_t + 1}\right) + 1$, trong đó $N$ là tổng số tài liệu và $n_t$ là số tài liệu chứa từ $t$.

Trọng số TF-IDF cuối cùng là tích của TF và IDF.

\textbf{Bước 3: Truy vấn}

\begin{enumerate}
    \item Tiền xử lý truy vấn: tách từ, loại bỏ từ dừng.
    
    \item Tìm ứng viên: Sử dụng chỉ mục ngược để tìm tất cả tài liệu chứa ít nhất một từ trong truy vấn.
    
    \item Tính độ tương đồng cosine: Với mỗi tài liệu ứng viên, tính tích vô hướng của vector truy vấn và vector tài liệu, chia cho tích của độ dài hai vector. Công thức: $similarity(q, d) = \frac{\vec{q} \cdot \vec{d}}{|\vec{q}| \times |\vec{d}|}$.
    
    \item Sắp xếp và trả về top-K tài liệu có độ tương đồng cao nhất.
\end{enumerate}

\textbf{Tối ưu hóa:} Thay vì lưu toàn bộ ma trận TF-IDF (222,981 x 50,000), hệ thống chỉ tính toán TF-IDF cho các tài liệu ứng viên được tìm thấy qua chỉ mục ngược, giảm đáng kể chi phí tính toán.


\subsection{Mô hình BM25}

BM25 được cài đặt dựa trên công thức chuẩn với các tham số $k_1 = 1.5$ và $b = 0.75$.

\textbf{Cấu trúc dữ liệu:}

Các giá trị cần tính trước được lưu trong chỉ mục:
\begin{itemize}
    \item $f(q_i, D)$: term frequency trong mỗi tài liệu
    \item $|D|$: độ dài mỗi tài liệu
    \item $N$: tổng số tài liệu
    \item $\text{avgdl}$: độ dài trung bình của corpus
\end{itemize}

\textbf{Thuật toán tính điểm BM25:}

\begin{algorithm}[H]
\caption{ScoreBM25}
\SetAlgoLined
\KwIn{$Q$ - truy vấn (danh sách các term), $D$ - tài liệu}
\KwOut{$score$ - điểm số BM25}
$score \leftarrow 0$\;
\ForEach{$q_i \in Q$}{
    \If{$q_i \notin D$}{
        \Continue\;
    }
    $f \leftarrow \text{getTF}(q_i, D)$\;
    $n_q \leftarrow \text{getDocFreq}(q_i)$\;
    $idf \leftarrow \ln\left(\frac{N - n_q + 0.5}{n_q + 0.5} + 1\right)$\;
    $numerator \leftarrow f \times (k_1 + 1)$\;
    $denominator \leftarrow f + k_1 \times \left(1 - b + b \times \frac{|D|}{\text{avgdl}}\right)$\;
    $score \leftarrow score + idf \times \frac{numerator}{denominator}$\;
}
\Return{$score$}
\end{algorithm}

\textbf{Thuật toán truy vấn BM25:}

\begin{algorithm}[H]
\caption{QueryBM25}
\SetAlgoLined
\KwIn{$Q$ - chuỗi truy vấn, $k$ - số kết quả cần trả về}
\KwOut{$R$ - danh sách top-$k$ tài liệu}
$Q \leftarrow \text{preprocess}(Q)$\;
$\text{candidates} \leftarrow \text{getCandidates}(Q)$\;
$\text{scores} \leftarrow []$\;
\ForEach{$D \in \text{candidates}$}{
    $s \leftarrow \text{ScoreBM25}(Q, D)$\;
    $\text{scores}.\text{append}((D, s))$\;
}
$\text{scores}.\text{sort}(\text{by score descending})$\;
$R \leftarrow \text{scores}[0:k]$\;
\Return{$R$}
\end{algorithm}

\textbf{So sánh với VSM:} BM25 có hai cải tiến chính so với VSM:
\begin{itemize}
    \item Độ bão hòa TF (TF saturation): Giới hạn đóng góp của tần suất cao.
    \item Chuẩn hóa độ dài (Length normalization): Cân bằng giữa tài liệu dài và ngắn.
\end{itemize}

Kết quả thực nghiệm cho thấy BM25 vượt trội hơn VSM đáng kể (MRR 0.5269 so với 0.3353).


\section{Kết quả và đánh giá}
Tiến hành đánh giá kết quả trên tập có 2000 truy vấn được lấy từ 2000 truy vấn đầu tiên khi kết hợp tập dữ liệu huấn luyện và tập kiểm thử.

Pipeline đánh giá được thực hiện như sau:
\begin{algorithm}
\caption{Pipeline đánh giá}
\textbf{Input:} search\_engine, ground\_truth: Dict[qid, List[cid]], queries: Dict[qid, question]

\begin{enumerate}
    \item \textbf{Với mỗi phương pháp:}\\
    Với mỗi method trong \{Boolean, VSM, BM25, ES\_Normal, ES\_Processed\}:
    \begin{enumerate}
        \item \textbf{Với mỗi truy vấn:}\\
        Với mỗi $(qid, question)$ trong queries:
        \begin{itemize}
            \item Thực hiện tìm kiếm: $results \gets search\_engine.search(question, method, top\_n=10)$
            \item Trích xuất ID: $retrieved\_ids$, $true\_ids$
            \item Tính các chỉ số: P@K, R@K, F1@K, RR, AP cho mỗi K
        \end{itemize}
    \end{enumerate}
    \item \textbf{Tổng hợp:} Tính trung bình của tất cả các chỉ số
\end{enumerate}
\end{algorithm}


\subsection{Kết quả đánh giá tổng quan}

Bảng \ref{tab:overall-results} trình bày kết quả đánh giá tổng quan của 5 phương pháp truy xuất trên các chỉ số Precision@1, Precision@5, Precision@10, Recall@10, F1@10, MRR và MAP.

\begin{table}[ht]
\centering
\begin{tabular}{l c c c c c c c}
\hline
Phương pháp & P@1 & P@5 & P@10 & R@10 & F1@10 & MRR & MAP \\
\hline
\textbf{ES (Đã xử lý)} & \textbf{0.4132} & \textbf{0.1349} & \textbf{0.0771} & \textbf{0.7708} & \textbf{0.1401} & \textbf{0.5284} & \textbf{0.5284} \\
\textbf{BM25} & \textbf{0.4120} & \textbf{0.1358} & \textbf{0.0770} & \textbf{0.7695} & \textbf{0.1399} & \textbf{0.5269} & \textbf{0.5269} \\
ES (Bình thường) & 0.4012 & 0.1339 & 0.0737 & 0.7372 & 0.1340 & 0.5136 & 0.5136 \\
VSM & 0.2195 & 0.0963 & 0.0627 & 0.6270 & 0.1140 & 0.3353 & 0.3353 \\
Boolean & 0.1305 & 0.0541 & 0.0306 & 0.3060 & 0.0556 & 0.1900 & 0.1900 \\
\hline
\end{tabular}
\caption{Kết quả đánh giá tổng quan các phương pháp}
\label{tab:overall-results}
\end{table}

Từ bảng kết quả, có thể thấy rằng Elasticsearch với chỉ mục đã xử lý đạt hiệu suất cao nhất với MRR và MAP đạt 0.5284. BM25 đứng thứ hai với MRR 0.5269, chỉ kém ES Processsed 0.0015. Điều này cho thấy thuật toán BM25 cài đặt thủ công có thể cạnh tranh được với Elasticsearch - một hệ thống tìm kiếm thương mại đã được tối ưu hóa.

ES với chỉ mục bình thường đạt MRR 0.5136, thấp hơn 0.0148 so với ES đã xử lý. Sự khác biệt này chứng minh tầm quan trọng của tiền xử lý văn bản tiếng Việt trước khi đưa vào hệ thống tìm kiếm. Khi văn bản được tách từ bằng underthesea, Elasticsearch có thể tìm kiếm chính xác hơn các từ ghép tiếng Việt.

VSM đạt MRR 0.3353, thấp hơn đáng kể so với BM25 và Elasticsearch. Nguyên nhân chính là VSM sử dụng TF-IDF đơn giản mà không có các cơ chế bão hòa TF và chuẩn hóa độ dài tài liệu như BM25.

Mô hình Boolean có hiệu suất thấp nhất với MRR chỉ 0.1900. Điều này là do mô hình Boolean không có cơ chế xếp hạng - tất cả tài liệu khớp đều được coi là có cùng mức độ liên quan, dẫn đến việc tài liệu liên quan có thể bị chìm trong hàng nghìn tài liệu khớp khác.

Các bảng chi tiết cho từng phương pháp:

% BM25
\begin{table}[H]
\centering
\begin{tabular}{c c c c}
\hline
K & Precision@K & Recall@K & F1@K \\
\hline
1  & 0.4120 & 0.4120 & 0.4120 \\
2  & 0.2685 & 0.5370 & 0.3580 \\
3  & 0.2033 & 0.6100 & 0.3050 \\
5  & 0.1358 & 0.6790 & 0.2263 \\
10 & 0.0769 & 0.7695 & 0.1399 \\
\hline
\end{tabular}
\caption{Kết quả đánh giá mô hình BM25}
\end{table}

% Boolean
\begin{table}[H]
\centering
\begin{tabular}{c c c c}
\hline
K & Precision@K & Recall@K & F1@K \\
\hline
1  & 0.1305 & 0.1305 & 0.1305 \\
2  & 0.0995 & 0.1990 & 0.1327 \\
3  & 0.0778 & 0.2335 & 0.1167 \\
5  & 0.0541 & 0.2705 & 0.0902 \\
10 & 0.0306 & 0.3060 & 0.0556 \\
\hline
\end{tabular}
\caption{Kết quả đánh giá mô hình Boolean}
\end{table}

% VSM
\begin{table}[H]
\centering
\begin{tabular}{c c c c}
\hline
K & Precision@K & Recall@K & F1@K \\
\hline
1  & 0.2195 & 0.2195 & 0.2195 \\
2  & 0.1633 & 0.3265 & 0.2177 \\
3  & 0.1323 & 0.3970 & 0.1985 \\
5  & 0.0963 & 0.4815 & 0.1605 \\
10 & 0.0627 & 0.6270 & 0.1140 \\
\hline
\end{tabular}
\caption{Kết quả đánh giá mô hình VSM}
\end{table}

% ES Processed
\begin{table}[H]
\centering
\begin{tabular}{rccc}
\hline
K & Precision@K & Recall@K & F1@K \\
\hline
1  & 0.4132 & 0.4132 & 0.4132 \\
2  & 0.2716 & 0.5432 & 0.3621 \\
3  & 0.2043 & 0.6128 & 0.3064 \\
5  & 0.1349 & 0.6744 & 0.2248 \\
10 & 0.0771 & 0.7708 & 0.1401 \\
\hline
\end{tabular}
\caption{Kết quả đánh giá ES với chỉ mục đã được xử lý}
\end{table}

% ES Normal
\begin{table}[H]
\centering
\begin{tabular}{rccc}
\hline
K & Precision@K & Recall@K & F1@K \\
\hline
1  & 0.4012 & 0.4012 & 0.4012 \\
2  & 0.2644 & 0.5288 & 0.3525 \\
3  & 0.1988 & 0.5964 & 0.2982 \\
5  & 0.1339 & 0.6696 & 0.2232 \\
10 & 0.0737 & 0.7372 & 0.1340 \\
\hline
\end{tabular}
\caption{Kết quả đánh giá ES với chỉ mục bình thường}
\end{table}


\subsection{Nhận xét kết quả}

Từ kết quả đánh giá, có thể rút ra một số nhận xét quan trọng sau đây:

\textbf{So sánh BM25 và ES đã xử lý:}
Hai phương pháp có hiệu suất gần như tương đương với MRR chênh lệch chỉ 0.0015. ES nhỉnh hơn một chút do thuật toán BM25 trong Elasticsearch đã được tối ưu hóa. Cả hai đều đạt xấp xỉ 77\% Recall@10, nghĩa là trong top 10 kết quả có chứa khoảng 77\% tài liệu liên quan - một kết quả rất tốt với tập dữ liệu có tỷ lệ positive thấp. Điều này chứng minh rằng thuật toán BM25 cài đặt thủ công hoàn toàn có thể cạnh tranh được với các hệ thống tìm kiếm thương mại. Sự chênh lệch nhỏ giữa hai phương pháp có thể đến từ việc Elasticsearch sử dụng nhiều tối ưu hóa bổ sung như caching, bloom filters, và các thuật toán phân tán.

\textbf{So sánh hai cách đánh chỉ mục của ES:}
ES với chỉ mục đã được xử lý tách từ cho kết quả tốt hơn. Cụ thể, MRR tăng từ 0.5136 lên 0.5284, tương đương với mức cải thiện 2.88\%. Điều này chứng tỏ tiền xử lý văn bản tiếng Việt giúp cải thiện hiệu suất tìm kiếm đáng kể. Khi truy vấn được tách từ (tokenize) trước bằng underthesea, Elasticsearch sẽ tìm chính xác hơn các từ ghép như "truy\_cứu", "trách\_nhiệm", "hình\_sự" thay vì tìm các token riêng lẻ không có nghĩa.

Sự khác biệt rõ rệt nhất ở Recall@10: ES đã xử lý đạt 77.08\% trong khi ES bình thường chỉ đạt 73.72\%. Điều này có nghĩa là với cùng 10 kết quả trả về, ES đã xử lý tìm được nhiều tài liệu liên quan hơn 3.36\%. Với tập dữ liệu mà mỗi truy vấn chỉ có 1-3 tài liệu liên quan trong hàng trăm nghìn tài liệu, sự cải thiện này là đáng kể.

\textbf{BM25 vs VSM:}
BM25 vượt trội hơn VSM một cách đáng kể với MRR 0.5269 so với 0.3353 (cải thiện 57.1\%). Nguyên nhân chính đến từ các đặc tính tốt hơn của BM25:
\begin{itemize}
    \item \textbf{Độ bão hòa TF (TF saturation):} Trong VSM, khi một từ xuất hiện nhiều lần trong tài liệu, điểm số tăng tuyến tính. Điều này có thể khiến các tài liệu dài với nhiều lần lặp lại của từ khóa bị đẩy lên cao dù không thực sự liên quan. BM25 giải quyết vấn đề này bằng hàm bão hòa, làm giảm ảnh hưởng của từ xuất hiện quá nhiều lần.
    \item \textbf{Tham số chuẩn hóa độ dài tài liệu ($b$):} BM25 sử dụng tham số $b$ để cân bằng giữa văn bản dài và ngắn. VSM không có cơ chế này, dẫn đến các tài liệu dài có xu hướng được ưu tiên hơn do có nhiều từ khớp hơn.
    \item \textbf{IDF bão hòa:} BM25 sử dụng công thức IDF có tính bão hòa, trong khi VSM sử dụng IDF log đơn giản. Điều này giúp BM25 xử lý tốt hơn các từ rất phổ biến trong corpus.
\end{itemize}

\textbf{Mô hình Boolean:}
Mô hình Boolean có hiệu suất thấp nhất trong các phương pháp được đánh giá. MRR chỉ đạt 0.1900, thấp hơn 3 lần so với BM25. Nguyên nhân chính là không có xếp hạng - tất cả tài liệu liên quan được coi là có cùng mức độ liên quan. Khi một truy vấn có nhiều từ (ví dụ: "phạm tội và trộm cắp và tài sản"), kết quả trả về có thể rất ít. Ngược lại, khi truy vấn ngắn (ví dụ: "luật"), kết quả trả về có thể rất nhiều (hàng nghìn tài liệu). Điều này dẫn đến việc tài liệu liên quan thực sự có thể bị chìm trong hàng nghìn tài liệu khác.

Precision@1 của Boolean chỉ đạt 13.05\%, có nghĩa là chỉ khoảng 1/8 truy vấn có tài liệu liên quan ở vị trí đầu tiên. Đây là hệ quả trực tiếp của việc không có xếp hạng - tài liệu liên quan có thể nằm ở bất kỳ vị trí nào trong danh sách kết quả.

\textbf{Xu hướng khi K thay đổi:}
Khi K tăng:
\begin{itemize}
    \item Precision@K giảm (vì thêm nhiều tài liệu không liên quan).
    \item Recall@K tăng (vì tìm được nhiều tài liệu liên quan hơn).
    \item Đây là đặc tính cơ bản của Precision-Recall trade-off.
\end{itemize}

Cụ thể, với BM25:
\begin{itemize}
    \item Precision@1 = 41.20\% → Precision@10 = 7.69\% (giảm 81.3\%)
    \item Recall@1 = 41.20\% → Recall@10 = 76.95\% (tăng 86.8\%)
\end{itemize}

Xu hướng này cho thấy nếu người dùng chỉ xem kết quả đầu tiên, cần hệ thống có Precision@1 cao. Nếu người dùng sẵn sàng xem nhiều kết quả, Recall@10 cao sẽ đảm bảo hầu hết tài liệu liên quan được hiển thị. Trong thực tế, người dùng thường xem 5-10 kết quả đầu tiên, do đó Recall@10 và MRR là các chỉ số quan trọng nhất.

\textbf{Phân tích lỗi:}
Mặc dù kết quả tổng thể khả quan, vẫn còn một số trường hợp hệ thống hoạt động chưa tốt:
\begin{itemize}
    \item \textbf{Truy vấn dài:} Khi truy vấn có nhiều từ, khả năng tìm được tài liệu khớp giảm do yêu cầu tất cả các từ đều phải xuất hiện.
    \item \textbf{Từ đồng nghĩa:} Hệ thống không xử lý từ đồng nghĩa, ví dụ "phạm tội" và "tội phạm" được coi là hai từ khác nhau.
    \item \textbf{Từ không có trong corpus:} Nếu truy vấn chứa từ không xuất hiện trong bất kỳ tài liệu nào, kết quả sẽ rất hạn chế.
\end{itemize}


\section{Giao diện người dùng (UI)}
Đề tài xây dựng giao diện người dùng để chạy thử nghiệm và trình bày kết quả. Giao diện được xây dựng dựa trên thư viện Streamlit \cite{streamlit}, một thư viện Python cho phép tạo giao diện web nhanh chóng từ các script. Streamlit đặc biệt phù hợp cho các dự án nghiên cứu và thử nghiệm vì nó cho phép tạo giao diện tương tác mà không cần viết code frontend riêng biệt.

Lý do chọn Streamlit cho đề tài:
\begin{itemize}
    \item Dễ dàng tích hợp với các model và xử lý dữ liệu Python.
    \item Không cần viết HTML/CSS/JavaScript.
    \item Hỗ trợ tự động tạo giao diện responsive.
    \item Có sẵn các component như sidebar, text input, select box, table, v.v.
    \item Cho phép hiển thị kết quả dưới dạng bảng, chart, văn bản.
\end{itemize}

Các tính năng của giao diện:
\begin{enumerate}
    \item \textbf{Chọn phương pháp tìm kiếm:} Sidebar cho phép người dùng chọn 1 trong 5 phương pháp:
    \begin{itemize}
        \item Boolean
        \item Mô hình không gian vector (VSM)
        \item BM25 (mặc định)
        \item Elasticsearch (Thông thường)
        \item Elasticsearch (Đã xử lý)
    \end{itemize}
    \item \textbf{Chọn số kết quả:} Ô chọn top N kết quả trả về (mặc định là 10).
    \item \textbf{Nhập truy vấn:} Ô nhập câu hỏi tìm kiếm với placeholder gợi ý.
    \item \textbf{Nút tìm kiếm:} Bắt đầu tìm kiếm khi người dùng nhập xong.
    \item \textbf{Hiển thị kết quả:} Bao gồm:
    \begin{itemize}
        \item Rank: Thứ tự (1, 2, ...)
        \item Document ID: ID của văn bản trong corpus
        \item Score: Điểm số (hoặc "N/A (Match)" cho Boolean)
        \item Nội dung: Văn bản gốc đầy đủ
    \end{itemize}
\end{enumerate}

Giao diện chương trình minh họa như trên hình \ref{fig:demo-ui}.

\begin{figure}[H]
    \centering
    \includegraphics[width=1\linewidth]{Hinhve/demo-ui.png}
    \caption{Giao diện chương trình minh họa kết quả}
    \label{fig:demo-ui}
\end{figure}

Quy trình xử lý khi người dùng thực hiện tìm kiếm:
\begin{enumerate}
    \item Người dùng nhập câu hỏi vào ô tìm kiếm.
    \item Nhấn nút "Tìm kiếm" hoặc nhấn Enter.
    \item Hệ thống tiền xử lý câu hỏi (tách từ, chuẩn hóa).
    \item Gọi phương pháp tìm kiếm đã chọn.
    \item Nhận kết quả và hiển thị lên giao diện.
\end{enumerate}

Một số lưu ý khi sử dụng giao diện:
\begin{itemize}
    \item Với phương pháp Boolean, người dùng có thể sử dụng các toán tử AND, OR, NOT (viết hoa).
    \item Với các phương pháp khác, người dùng nhập câu hỏi bình thường.
    \item Thời gian tìm kiếm phụ thuộc vào phương pháp và độ phức tạp của truy vấn.
    \item Kết quả được sắp xếp theo điểm số (ngoại trừ Boolean).
\end{itemize}

\end{document}
